% TEMPLATE-MILC-v3.tex - plantilla maestra UNICA de las guias MILC v3.
%
% La usan las tres familias de guias, cada una con su driver:
%   · Grados 6-11  -> scripts/build-guias-g11.py
%   · Bebras       -> scripts/build-guias-bebras.py
%   · Semillero    -> scripts/build-guias-semillero.py
%
% Antes habia tres copias casi identicas (~400 lineas iguales, 6 distintas):
% cada mejora habia que aplicarla tres veces, y en la practica no se hacia
% (el motor de recursos vivia solo en dos de ellas). Lo que varia entre
% familias esta parametrizado en 6 placeholders que cada driver llena
% (se nombran aqui SIN su sintaxis real: el builder reemplaza por texto plano,
% tambien dentro de los comentarios, y un valor multilinea rompe el preambulo):
%
%   HEADER_IZQ        encabezado izquierdo de pagina
%   HEADER_DER        encabezado derecho de pagina
%   PORTADA_ETIQUETA  rotulo sobre el numero grande (GRADO/MOMENTO/MODULO)
%   PORTADA_SERIE     linea de serie de la portada
%   PORTADA_ID        identificador de la guia en la portada
%   PORTADA_CIRCULO   el circulito de grado (°), o vacio si no aplica
%
% El resto de claves las documenta content/guias/_SCHEMA.md. El script valida
% que no queden placeholders sin reemplazar antes de escribir el archivo final.

\documentclass[11pt,letterpaper]{article}
\usepackage[margin=2.0cm]{geometry}
\usepackage[table]{xcolor}
\usepackage{fontspec}
\usepackage[spanish]{babel}
\usepackage{tikz}
% Librerías TikZ para los diagramas de las guías (bloque `recursos:`):
%  · babel  → babel en español vuelve activos caracteres como «>», y TikZ los usa
%             en sus opciones (>=stealth, ->). Sin esta librería, cualquier
%             diagrama con flechas revienta con «\language@active@arg> has an
%             extra }». Debe ir ANTES de que se dibuje nada.
%  · positioning        → sintaxis «below left=2pt and 3pt».
%  · decorations.pathreplacing → llaves ({brace}) para señalar rangos.
\usetikzlibrary{babel,positioning,decorations.pathreplacing}
\usepackage{graphicx}
% Circuitos: el trabajo de electrónica se hace hoy con micro:bit y sus sensores
% integrados (MakeCode, sin cableado), así que no se necesitan esquemáticos.
% Si algún día entran montajes con protoboard, basta descomentar esta línea:
% los diagramas .tex/.tikz declarados en `recursos:` ya se incrustan solos.
% \usepackage{circuitikz}
\usepackage[most]{tcolorbox}
\usepackage{tabularx}
\usepackage{array}
\usepackage{fancyhdr}
\usepackage{titlesec}
\usepackage[document]{ragged2e}  % bandera a la derecha en todo el cuerpo
\usepackage{enumitem}
\usepackage{microtype}

% Helvetica Neue no trae la flecha U+2192, asi que cada «→» escrito literalmente
% en el YAML se compilaba a NADA, en silencio (462 flechas en 61 guias: rutas del
% tipo «paleta → tipografia → lockup» salian rotas). Mapeamos el caracter a su
% version matematica, que si tiene glifo. Asi el autor puede seguir escribiendo
% «→» comodamente en el YAML.
\usepackage{newunicodechar}
\newunicodechar{→}{\ensuremath{\rightarrow}}
% Mismo problema con los simbolos de marca y vinetas: el «✓» de «una palomita ✓»
% salia como cuadrito vacio en el PDF de todas las guias que lo usaban.
\usepackage{pifont}
\newunicodechar{✓}{\ding{51}}
\newunicodechar{✗}{\ding{55}}
\newunicodechar{✔}{\ding{52}}
\newunicodechar{•}{\textbullet}
\newunicodechar{≥}{\ensuremath{\geq}}
\newunicodechar{≤}{\ensuremath{\leq}}

\setmainfont{Helvetica Neue}
\setsansfont{Helvetica Neue}
\setlength{\parindent}{0pt}
\setlength{\parskip}{6pt}
% Sin particion de palabras: en 6.º-7.º una palabra cortada por guion al final
% de linea («Comuni-cacion») frena la lectura mucho mas que una linea corta.
\hyphenpenalty=10000
\exhyphenpenalty=10000
\pretolerance=2000
\tolerance=3000
\setlength{\emergencystretch}{3em}
\linespread{1.10}
\renewcommand{\arraystretch}{1.25}
\setlist{leftmargin=5mm,itemsep=1mm,topsep=1mm}
\newcolumntype{Y}{>{\RaggedRight\arraybackslash}X}

\definecolor{milcMagenta}{HTML}{E6007E}
\definecolor{milcVino}{HTML}{5A0038}
\definecolor{milcTurquesa}{HTML}{0093A5}
\definecolor{milcVerde}{HTML}{58A923}
\definecolor{milcMostaza}{HTML}{E5B400}
\definecolor{milcOcre}{HTML}{B86600}
\definecolor{milcNegro}{HTML}{241816}
\definecolor{milcGris}{HTML}{F7F7F4}
\definecolor{milcLinea}{HTML}{E8E2DD}
\definecolor{milcGradePrimary}{HTML}{FF2D87}
\definecolor{milcGradeDark}{HTML}{9F1B56}
\definecolor{milcGradeSoft}{HTML}{FFE0EE}
\definecolor{milcGradeAccent}{HTML}{CFE7FF}
\definecolor{milcGradeText}{HTML}{FFFFFF}
\color{milcNegro}

\pagestyle{fancy}
\fancyhf{}
\lhead{\small Tecnología e Informática · Grado 10}
\rhead{\small Guía 17 · MILC v3}
\cfoot{\small\thepage}
\renewcommand{\headrulewidth}{0pt}

\newtcolorbox{titlebox}[1]{
  enhanced,colback=#1,colframe=#1,arc=7mm,boxrule=0pt,
  left=7mm,right=7mm,top=3mm,bottom=3mm,
  before skip=8pt,after skip=8pt,
  fontupper=\sffamily\bfseries\Large\color{white}
}

\newtcolorbox{softbox}[3]{
  enhanced,colback=#2,colframe=#1,borderline west={4pt}{0pt}{#1},
  arc=2mm,boxrule=.4pt,left=4mm,right=4mm,top=3mm,bottom=3mm,
  before skip=6pt,after skip=8pt,title={#3},coltitle=white,
  colbacktitle=#1,fonttitle=\sffamily\bfseries,fontupper=\RaggedRight,
  attach boxed title to top left={xshift=3mm,yshift=-2mm},
  boxed title style={arc=2mm, boxrule=0pt}
}

\newtcolorbox{drawbox}[2]{
  enhanced,colback=white,colframe=#1,arc=4mm,boxrule=1pt,
  left=5mm,right=5mm,top=4mm,bottom=4mm,before skip=8pt,after skip=8pt,
  title={#2},coltitle=white,colbacktitle=#1,fonttitle=\sffamily\bfseries,
  fontupper=\RaggedRight,
  attach boxed title to top left={xshift=4mm,yshift=-2mm},
  boxed title style={arc=3mm, boxrule=0pt}
}

\newtcolorbox{infoband}[1]{
  enhanced,colback=#1!8,colframe=#1!50,arc=2mm,boxrule=.5pt,
  left=4mm,right=4mm,top=2mm,bottom=2mm,before skip=4pt,after skip=4pt,
  fontupper=\small\RaggedRight
}

% Puente narrativo entre secciones (contrato editorial Conéctate).
% Da continuidad: "Acabas de X. Ahora vas a Y porque Z."
% Visualmente: banda izquierda en color del grado, texto en itálica pequeña.
\newtcolorbox{puentebox}{
  enhanced,colback=milcGradeSoft,colframe=milcGradeDark,
  borderline west={3pt}{0pt}{milcGradeDark},
  arc=0mm,boxrule=0pt,left=5mm,right=4mm,top=2.5mm,bottom=2.5mm,
  before skip=4pt,after skip=8pt,
  fontupper=\itshape\small\color{milcNegro}\RaggedRight
}
% La flecha va en modo matematico a proposito: Helvetica Neue NO tiene el glifo
% U+2192, asi que \textrightarrow se compilaba a NADA («Missing character: There
% is no → in font Helvetica Neue») y el puente salia sin flecha. En modo
% matematico la flecha viene de la fuente matematica y si se dibuja.
\newcommand{\puente}[1]{%
  \begin{puentebox}%
    \textcolor{milcGradeDark}{$\rightarrow$}\hspace{2mm}#1%
  \end{puentebox}%
}

\newcommand{\linead}[1]{\noindent\color{milcLinea}\rule{#1}{0.5pt}\color{milcNegro}}
\newcommand{\respuesta}[1]{\par\vspace{2mm}\foreach \n in {1,...,#1}{\noindent\color{milcLinea}\rule{\linewidth}{0.5pt}\par\vspace{5mm}}\color{milcNegro}}
\newcommand{\checkbox}{\textcolor{milcGradeDark}{\fbox{\phantom{x}}}\hspace{5pt}}

\newcommand{\phasecard}[4]{
\begin{tcolorbox}[enhanced,colback=#3,colframe=#2,arc=3mm,boxrule=.5pt,height=3.05cm,valign=center,halign=center,left=1.5mm,right=1.5mm,top=2mm,bottom=2mm]
{\sffamily\bfseries\fontsize{22}{24}\selectfont\textcolor{#2}{#1}}\par
{\sffamily\bfseries\small #4}
\end{tcolorbox}
}

% ── Piezas del rediseno 2026 (legibilidad 6.º-11.º) ───────────────────
% Banda de estacion: reemplaza el titlebox macizo. Titulo a la izquierda,
% dato practico (tiempo/modalidad) a la derecha, en una sola linea.
\newcommand{\milcestacion}[2]{%
\begin{tcolorbox}[enhanced,colback=#1,colframe=#1,arc=2mm,boxrule=0pt,
  left=5mm,right=5mm,top=2.6mm,bottom=2.6mm,before skip=11pt,after skip=7pt]
{\sffamily\bfseries\fontsize{15}{18}\selectfont\color{white}#2}
\end{tcolorbox}}

% Tarjeta de paso numerada: sustituye las tablas «Pilar / Aplicacion», que
% eran formato de consulta docente, no de estudiante. El numero va en un
% circulo sobre el borde para que la secuencia se lea de un vistazo.
\newcommand{\milcpaso}[3]{%
\begin{tcolorbox}[enhanced,colback=white,colframe=milcLinea,arc=1.5mm,boxrule=.9pt,
  left=14mm,right=4mm,top=3mm,bottom=3mm,before skip=4pt,after skip=4pt,
  fontupper=\RaggedRight,
  overlay={\node[anchor=north west,circle,fill=#2,text=white,inner sep=0pt,
    minimum size=7.4mm,font=\sffamily\bfseries\small]
    at ([xshift=3.6mm,yshift=-3.2mm]frame.north west) {#1};}]
#3
\end{tcolorbox}}

% Casilla marcable de tamano util con lapiz (la anterior era \fbox{\phantom{x}},
% de alto variable segun el texto que la seguia).
\newcommand{\milcmarca}{\framebox[3.8mm]{\rule{0pt}{3.4mm}}}

% Fila marcable de lista suelta (checks de la estacion 1).
\newcommand{\milccheckrow}[1]{%
\par\vspace{1.8mm}\noindent
\makebox[6.4mm][l]{\raisebox{-.55mm}{\textcolor{milcNegro}{\milcmarca}}}%
\parbox[t]{\dimexpr\linewidth-6.4mm\relax}{\RaggedRight #1}\par}

% Celda marcable dentro de la rubrica (tabla de autoevaluacion). Misma
% sangria francesa, pero pensada para vivir dentro de una columna X.
\newcommand{\milcopcion}[1]{%
\makebox[5.2mm][l]{\raisebox{-.55mm}{\textcolor{milcNegro}{\milcmarca}}}%
\parbox[t]{\dimexpr\linewidth-5.2mm\relax}{\RaggedRight #1}}

% ── Recursos visuales de la guía ──────────────────────────────────────
% Declarados en el bloque `recursos:` del YAML y emitidos por el builder.
% \guiaFigura   → imagen rasterizada o PDF (foto, carta celeste, montaje).
% \guiaDiagrama → fuente TikZ/circuitikz (.tex) incrustada como vector:
%                 es la vía para circuitos y esquemáticos editables.
% #1 = ruta relativa al .tex · #2 = pie (puede ir vacío)
\newcommand{\guiaFigura}[2]{%
\begin{center}
\includegraphics[width=0.88\linewidth,height=0.38\textheight,keepaspectratio]{#1}\\[1.5mm]
{\footnotesize\itshape\textcolor{milcNegro!75}{#2}}
\end{center}
}

\newcommand{\guiaDiagrama}[2]{%
\begin{center}
\input{#1}\\[1.5mm]
{\footnotesize\itshape\textcolor{milcNegro!75}{#2}}
\end{center}
}

\begin{document}
\thispagestyle{empty}

% ── Portada ───────────────────────────────────────────────────────────
% Antes: un rectangulo de color a pagina completa. Se veia bien en pantalla,
% pero en la fotocopiadora del colegio se comia el toner y salia gris sucio.
% Ahora la banda ocupa el tercio superior y el resto va en blanco.
\begin{tikzpicture}[remember picture,overlay]
  \path[fill=milcGradePrimary]
    (current page.north west) rectangle ([yshift=-10.6cm]current page.north east);

  \node[anchor=north west,text=milcGradeText] at ([xshift=2.0cm,yshift=-1.55cm]current page.north west)
    {\sffamily\bfseries\fontsize{12}{14}\selectfont GRADO};
  \node[anchor=north west,text=milcGradeText,opacity=.85] at ([xshift=2.0cm,yshift=-2.35cm]current page.north west)
    {\sffamily\bfseries\fontsize{8.5}{10}\selectfont SERIE GUÍAS · TECNOLOGÍA E INFORMÁTICA};
  \node[anchor=north east,text=milcGradeText,opacity=.85,align=right] at ([xshift=-2.0cm,yshift=-1.55cm]current page.north east)
    {\sffamily\bfseries\fontsize{9}{12}\selectfont I.E. Sor María Juliana\\Cartago, Valle del Cauca};

  \node[anchor=west,text=milcGradeText] (gradonum) at ([xshift=1.85cm,yshift=-7.1cm]current page.north west)
    {\sffamily\bfseries\fontsize{112}{112}\selectfont 10};
  % El circulito de grado (°) solo aplica a las guias de grado («11°»); en
  % Bebras y Semillero el numero grande es un momento/modulo, no un grado.
  % Se posiciona relativo al nodo `gradonum`, no en coordenadas absolutas.
  \draw[milcGradeText,line width=1.6pt]
    ([xshift=2.0mm,yshift=-4.5mm]gradonum.north east) circle (.15cm);

  \node[anchor=west,text=milcGradeText,text width=11.4cm,align=left] at ([xshift=1.5cm,yshift=0.5cm]gradonum.east)
    {
     {\sffamily\bfseries\fontsize{9}{11}\selectfont\MakeUppercase{Período 2 · Informes técnicos y comunicación profesional con IA}}\\[3mm]
     {\sffamily\bfseries\fontsize{21}{25}\selectfont\setlength{\baselineskip}{27pt}Informe comercial ---\\[4pt]propuesta concreta\\[4pt]para audiencia real}\\[3.5mm]
     {\sffamily\bfseries\fontsize{15}{18}\selectfont Guía 17-10-TIC}
    };
\end{tikzpicture}

\vspace*{9.4cm}
\vfill

{\sffamily\bfseries\fontsize{8}{10}\selectfont\textcolor{milcGradeDark}{LO QUE TE LLEVAS HOY}}

\begin{tcolorbox}[enhanced,colback=white,colframe=milcNegro,arc=2mm,boxrule=1.6pt,
  left=5mm,right=5mm,top=4mm,bottom=4mm,before skip=3pt,after skip=12pt,
  fontupper=\RaggedRight\fontsize{11}{15.5}\selectfont]
Informe comercial completo dirigido a una audiencia real concreta (consejo estudiantil, rector, coordinación académica, organización del barrio, microempresario de la familia) sobre propuesta de mejora o iniciativa real con (a) resumen ejecutivo de 1 página, (b) cuerpo argumentativo de 2-3 páginas con datos del análisis (sesión 6), (c) recomendaciones accionables con responsables sugeridos, (d) presupuesto estimado si aplica + identificación clara del destinatario y plan de entrega
\end{tcolorbox}

\begin{tcolorbox}[enhanced,colback=milcGris,colframe=milcGris,arc=2mm,boxrule=0pt,
  left=5mm,right=5mm,top=3.5mm,bottom=3.5mm,after skip=12pt]
{\sffamily\bfseries\fontsize{8}{10}\selectfont\textcolor{milcNegro!55}{ESTA GUÍA ES DE}}\\[3mm]
\begin{tabularx}{\linewidth}{X p{3.2cm} p{3.2cm}}
\linead{\linewidth} & \linead{3.0cm} & \linead{3.0cm}\\[-1mm]
{\fontsize{8.5}{10}\selectfont\textcolor{milcNegro!55}{Nombre completo}} &
{\fontsize{8.5}{10}\selectfont\textcolor{milcNegro!55}{Grupo}} &
{\fontsize{8.5}{10}\selectfont\textcolor{milcNegro!55}{Fecha}}\\
\end{tabularx}
\end{tcolorbox}

\vfill

\noindent\textcolor{milcLinea}{\rule{\linewidth}{1pt}}\\[2mm]
{\sffamily\bfseries\fontsize{9.5}{12}\selectfont Autor: Dr. Álvaro Cárdenas Orozco}\\[1mm]
{\sffamily\fontsize{8.5}{11}\selectfont\textcolor{milcNegro!55}{Metodología MILC · apertura ancestral · pensamiento computacional · triángulo Dussel–estoicismo–Floridi}}

\newpage

\section*{Apertura: Informe comercial --- para audiencia real}

\begin{softbox}{milcGradeDark}{milcGradeSoft}{Saber ancestral}
En las décadas del 60, 70 y 80 del siglo XX, cuando los bancos del centro de Cartago todavía atendían a los pequeños comerciantes del barrio, había una práctica que cualquier microempresario experimentado conocía: \textbf{cuando el tendero iba al banco a pedir un préstamo, no llegaba con palabras}. Llegaba con \textbf{un papel}. En ese papel venían: (a) las cuentas claras de la tienda los últimos 6 meses (ingresos y egresos por semana), (b) la propuesta concreta de qué iba a hacer con la plata pedida (comprar nevera, ampliar la tienda, traer mercancía nueva), (c) la cantidad exacta pedida en pesos, (d) cómo y en cuánto tiempo iba a pagar el préstamo. El gerente del banco recibía el papel, lo leía en 5 minutos, y decidía. Si el papel estaba bien hecho, el préstamo se aprobaba. Si era confuso o incompleto, el tendero volvía a su casa sin la plata. La sabiduría es ancestral y precisa: \textbf{frente a un tomador de decisiones, lo escrito vale más que lo hablado}. Ese papel que el tendero llevaba al banco era \textbf{informe comercial avant la lettre}: propuesta concreta dirigida a audiencia real con poder de decidir.
\end{softbox}

\begin{softbox}{milcTurquesa}{milcGris}{Saber contemporáneo}
Un \textbf{informe comercial} es texto profesional escrito para \textbf{persuadir a una audiencia específica con poder de decidir}. Se diferencia del informe técnico (sesión 2) en que su propósito principal es \emph{lograr que algo pase}: que aprueben un proyecto, que financien una iniciativa, que cambien una norma, que contraten un servicio. Sus 4 componentes obligatorios son: (1) \textbf{Resumen ejecutivo} (1 página): qué propones, a quién se dirige, qué se necesita decidir, qué beneficios trae. La regla profesional: el lector que solo lee el resumen ejecutivo debe poder decidir. (2) \textbf{Cuerpo argumentativo} (2-3 páginas): contexto, datos que sustentan, comparaciones, argumentos para superar objeciones probables. (3) \textbf{Recomendaciones accionables}: qué se propone hacer, quién es responsable, en qué plazo, con qué recursos. (4) \textbf{Presupuesto estimado} (si aplica): cuánto cuesta, de dónde sale el dinero, qué se obtiene. La regla profesional: \textbf{informe comercial sin pedido claro es informe perdido}. Debe terminar con \emph{``por lo tanto, solicito decisión sobre X''}.
\end{softbox}

\begin{softbox}{milcMostaza}{milcGris}{Pregunta-puente}
\textit{¿Qué sabía el tendero al llevar el papel con cuentas claras al banco, que el novato olvida cuando hace propuesta verbal sin documento? ¿Y por qué el resumen ejecutivo de 1 página es la pieza más importante del informe comercial?}
\end{softbox}

\begin{softbox}{milcVerde}{milcGris}{Saber y saber hacer}
Al terminar podrás: (1) \textbf{analizar} qué audiencia real con poder de decidir podría recibir tu informe comercial; (2) \textbf{aplicar} la estructura de 4 componentes a una propuesta concreta del colegio o barrio; (3) \textbf{crear} el informe comercial completo de 3-4 páginas con resumen ejecutivo + cuerpo + recomendaciones + presupuesto si aplica; (4) \textbf{evaluar} si tu informe lograría decisión positiva si llegara a la audiencia real.
\end{softbox}



\small
\begin{tabularx}{\textwidth}{>{\bfseries}p{3.0cm}Y}
\rowcolor{milcGradePrimary!12}Autor & Dr. Álvaro Cárdenas Orozco\\
DBA & Comunicación profesional con asistencia de IA y herramientas ofimáticas gratuitas (MEN, T\&I)\\
Referentes & Markdown como estándar abierto · Google Workspace · Estoicismo (claridad)\\
Producto final & Informe comercial completo dirigido a una audiencia real concreta (consejo estudiantil, rector, coordinación académica, organización del barrio, microempresario de la familia) sobre propuesta de mejora o iniciativa real con (a) resumen ejecutivo de 1 página, (b) cuerpo argumentativo de 2-3 páginas con datos del análisis (sesión 6), (c) recomendaciones accionables con responsables sugeridos, (d) presupuesto estimado si aplica + identificación clara del destinatario y plan de entrega\\
\end{tabularx}
\normalsize

\puente{Antes de empezar, mira el viaje completo. En la sesión 2 aprendiste informe técnico (informar); en la 6, análisis con IA. Hoy aprendes informe comercial (persuadir para decidir). Las próximas sesiones (correo profesional, mini-proyecto, sustentación) integrarán todo en propuestas reales.}

\milcestacion{milcGradeDark}{Ruta MILC: así vamos a aprender}
\begin{tabularx}{\textwidth}{>{\centering\arraybackslash}X>{\centering\arraybackslash}X>{\centering\arraybackslash}X>{\centering\arraybackslash}X}
\phasecard{1}{milcVerde}{milcGris}{Escucha\\[-1mm]\small Observo, escucho y pregunto.} &
\phasecard{2}{milcTurquesa}{milcGris}{Entender\\[-1mm]\small Sistematizo y ordeno.} &
\phasecard{3}{milcMagenta}{milcGris}{Hacer\\[-1mm]\small Construyo y aplico.} &
\phasecard{4}{milcVino}{milcGris}{Cierre\\[-1mm]\small Evalúo y reflexiono.}
\end{tabularx}


\puente{Antes de teorizar, vas a hacer una \emph{lluvia de propuestas}: anota 5 propuestas reales que harías a alguien con poder de decidir (consejo estudiantil, rector, JAC, microempresario familiar, alcaldía local). Cada una con su destinatario claro.}

\milcestacion{milcVerde}{Estación 1 de 3 · Escucha}

\begin{softbox}{milcVerde}{milcGris}{Escena para escuchar}
\textbf{Actividad 1 · ANALIZA --- Lluvia de 5 propuestas a audiencia real} (15 min · individual). Anota 5 propuestas concretas que harías a alguien con poder de decidir en tu entorno cercano. Cada una con destinatario nombrado. Ejemplos: (1) Al consejo estudiantil: propuesta de mejora del recreo. (2) A la coordinación académica: propuesta de revisión del horario. (3) A la JAC del barrio: propuesta de campaña ambiental. (4) Al microempresario familiar (tendero, vendedor): propuesta de digitalización con IA. (5) A la coordinación de informática: propuesta de mejora del laboratorio. Para cada una, anota qué decidiría el destinatario si aprobara.
\end{softbox}

{\sffamily\bfseries\fontsize{8}{10}\selectfont\textcolor{milcNegro!55}{MARCA CUANDO LO HAYAS HECHO}}
\milccheckrow{Anoté 5 propuestas reales con destinatario nombrado.}
\milccheckrow{Cada propuesta tiene decisión clara que se solicitaría.}
\milccheckrow{Reconozco que las 5 podrían ser informe comercial.}

\begin{infoband}{milcVerde}


\textbf{Lo que va al cuaderno (Actividad 1 · ANALIZA).} Título: «Actividad 1 --- 5 propuestas reales». \textbf{Formato:} lista de 5 con destinatario + decisión solicitada. \textbf{Sabes que terminaste cuando} eliges 1 con destinatario contactable.
\end{infoband}

\puente{Ya tienes 5 candidatos. Ahora vas a entender la \textbf{anatomía profesional del informe comercial}: 4 componentes, los 5 errores típicos que evitan la decisión, la diferencia con informe técnico.}

\milcestacion{milcTurquesa}{Estación 2 de 3 · Entender}

\begin{softbox}{milcTurquesa}{milcGris}{Lo que vas a entender}
Tus 5 propuestas tienen potencial. Ahora necesitas la \textbf{anatomía profesional} del informe comercial.
\end{softbox}

\milcpaso{1}{milcTurquesa}{Los \textbf{4 componentes obligatorios} se descomponen así: (1) \textbf{Resumen ejecutivo} (1 página): qué se propone (1-2 frases), para quién (destinatario nombrado), qué se necesita decidir (verbo de acción), beneficios principales (3-4 puntos), costos aproximados si aplica. Es la pieza más leída. (2) \textbf{Cuerpo argumentativo} (2-3 páginas): contexto del problema, datos que sustentan la necesidad, análisis (puede usar los hallazgos de la sesión 6), comparación con alternativas, argumentos para superar objeciones probables del destinatario. (3) \textbf{Recomendaciones accionables}: qué se propone hacer en pasos concretos, quién es responsable de cada paso, en qué plazo. (4) \textbf{Presupuesto estimado}: cuánto cuesta, de dónde sale el dinero, qué se obtiene a cambio (retorno cuantificable cuando sea posible).}
\milcpaso{2}{milcTurquesa}{La \textbf{diferencia con informe técnico} es crucial: el informe técnico informa para que alguien sepa; el comercial persuade para que alguien decida. Tono más cálido pero profesional. Beneficios destacados. Estructura orientada al cierre. La phronesis del oficio comercial consiste en \textbf{no esconder el pedido en el último párrafo}: el resumen ejecutivo ya debe decirlo.}
\milcpaso{3}{milcTurquesa}{Los \textbf{5 errores típicos} del informe comercial novato: (a) \textbf{Resumen ejecutivo vago}: \emph{``este es un informe sobre el recreo''} no es propuesta. (b) \textbf{Sin destinatario claro}: el informe es \emph{a quien corresponda} en lugar de \emph{a Carlos Pérez, coordinador}. (c) \textbf{Argumentos sin datos}: opiniones disfrazadas de propuesta. (d) \textbf{Recomendaciones sin responsables}: \emph{``alguien debería hacer X''}. (e) \textbf{Sin presupuesto cuando se necesita}: pedir cambios estructurales sin estimar costos.}
\milcpaso{4}{milcTurquesa}{Algoritmo: (1) elegir 1 propuesta de los 5 candidatos. (2) identificar destinatario real con nombre y cargo. (3) escribir resumen ejecutivo (1 página) que ya diga el pedido. (4) escribir cuerpo argumentativo con datos. (5) escribir recomendaciones con responsables y plazos. (6) calcular presupuesto si aplica. (7) preparar plan de entrega (cómo hacer llegar el informe).}

\begin{drawbox}{milcTurquesa}{Anatomía del informe comercial}
\textbf{Resumen ejecutivo:} 1 página con pedido claro. \\[1mm] \textbf{Cuerpo argumentativo:} 2-3 páginas con datos. \\[1mm] \textbf{Recomendaciones:} pasos + responsables + plazos. \\[1mm] \textbf{Presupuesto:} costos + financiamiento + retorno. \\[1mm] \textbf{Plan de entrega:} cómo llegar al destinatario.
\end{drawbox}

\begin{softbox}{milcMostaza}{milcGris}{Errores comunes y su corrección}
(1) Resumen ejecutivo vago sin pedido claro. (2) Sin destinatario nombrado. (3) Argumentos sin datos sustentables. (4) Recomendaciones sin responsables. (5) Sin presupuesto cuando se necesita. (6) Tono frío como informe técnico. (7) Pedido escondido al final.
\end{softbox}

\begin{infoband}{milcTurquesa}


\textbf{Lo que va al cuaderno (Actividad 2 · APLICA).} Título: «Actividad 2 --- Anatomía informe comercial». \textbf{Formato:} (a) los 4 componentes; (b) diferencia con informe técnico; (c) los 7 errores con marca. \textbf{Sabes que terminaste cuando} podrías auditar un informe comercial real.
\end{infoband}

\puente{Tienes el método. Toca aplicarlo: eliges 1 propuesta de los 5, identificas destinatario real, escribes los 4 componentes, preparas plan de entrega.}

\milcestacion{milcMagenta}{Estación 3 de 3 · Hacer}

\begin{softbox}{milcMagenta}{milcGris}{Cómo vas a construir tu producto}
\textbf{Actividad 3 · CREA + EVALÚA --- Informe comercial completo + plan de entrega} (40 min · individual). Hora de aplicar. Eliges propuesta, identificas destinatario, escribes los 4 componentes, preparas plan de entrega.
\end{softbox}

\milcpaso{1}{milcMagenta}{Tu informe comercial se construye en 7 pasos: (1) elegir 1 propuesta entre los 5 candidatos. (2) identificar destinatario real con nombre, cargo, contexto. (3) escribir resumen ejecutivo de 1 página que ya contenga el pedido. (4) escribir cuerpo argumentativo de 2-3 páginas con datos (puedes usar los hallazgos de la encuesta de sesión 5 y análisis de sesión 6). (5) escribir recomendaciones con responsables y plazos. (6) calcular presupuesto si aplica. (7) preparar plan de entrega (correo, reunión, carta presentada en persona).}
\milcpaso{2}{milcMagenta}{Sigue el patrón \textbf{``pedido en la primera página''}: el resumen ejecutivo no es introducción ni teaser; es el documento completo en miniatura. Quien lee solo el resumen ejecutivo debe poder decidir. La phronesis del informe comercial consiste en \textbf{no esconder el pedido}.}
\milcpaso{3}{milcMagenta}{El plan de entrega es pieza pequeña pero crucial: ¿cómo va a llegar tu informe al destinatario? Opciones: (a) entregar impreso en persona después de pedir cita. (b) enviar por correo electrónico con mensaje breve. (c) presentar en reunión con copia para llevar. (d) compartir vía la coordinación. Cada opción tiene protocolo distinto. La phronesis profesional incluye \textbf{pensar la entrega tan cuidadosamente como el contenido}.}
\milcpaso{4}{milcMagenta}{Algoritmo: (1) propuesta. (2) destinatario. (3) resumen ejecutivo. (4) cuerpo argumentativo. (5) recomendaciones. (6) presupuesto. (7) plan de entrega. (8) revisión final: ¿el destinatario podría decidir con base en este documento?}

\begin{drawbox}{milcMagenta}{Checklist antes de entregar}
\checkbox Propuesta elegida y destinatario nombrado. \\[1mm] \checkbox Resumen ejecutivo de 1 página con pedido claro. \\[1mm] \checkbox Cuerpo argumentativo con datos. \\[1mm] \checkbox Recomendaciones con responsables y plazos. \\[1mm] \checkbox Presupuesto si aplica. \\[1mm] \checkbox Tono cálido pero profesional. \\[1mm] \checkbox Plan de entrega preparado.
\end{drawbox}

\begin{softbox}{milcMagenta}{milcGris}{Plantilla de trabajo}
\textbf{Destinatario.} \textit{Nombre + cargo + contexto.} \\[1mm] \textbf{Resumen ejecutivo (1 pp).} \textit{Propuesta + pedido + beneficios + costos.} \\[1mm] \textbf{Cuerpo (2-3 pp).} \textit{Contexto + datos + argumentos.} \\[1mm] \textbf{Recomendaciones.} \textit{Pasos + responsables + plazos.} \\[1mm] \textbf{Presupuesto.} \textit{Costos + financiamiento + retorno.} \\[1mm] \textbf{Plan de entrega.} \textit{Cómo llegará al destinatario.}
\end{softbox}

\begin{infoband}{milcMagenta}


\textbf{Lo que va al cuaderno (Actividad 3 · CREA + EVALÚA).} Título: «Actividad 3 --- Informe comercial». \textbf{Formato:} (a) referencia al documento; (b) plan de entrega; (c) autoevaluación de si lograría decisión. \textbf{Sabes que terminaste cuando} podrías entregar el informe al destinatario esta semana.
\end{infoband}

\puente{Lo que sigue resume \emph{qué} entregas hoy: informe comercial completo + plan de entrega. El criterio principal: que el destinatario, recibiendo tu informe, pueda decidir en favor o en contra con base en lo que lee.}

\milcestacion{milcMostaza}{Producto de la sesión}
\textbf{Informe comercial completo + plan de entrega a destinatario real}.
\begin{softbox}{milcMostaza}{milcGris}{Criterios de calidad}
(1) Destinatario nombrado con cargo. (2) Resumen ejecutivo de 1 página con pedido. (3) Cuerpo argumentativo con datos. (4) Recomendaciones con responsables. (5) Presupuesto si aplica. (6) Plan de entrega preparado.
\end{softbox}

\puente{Antes de cerrar, mira el informe comercial desde las cinco dimensiones humanas. Persuadir con datos es respeto por el tomador de decisiones; pedir sin sustento es pérdida de tiempo de todos.}

\milcestacion{milcVino}{Evaluación liberadora · cinco dimensiones}

\begin{softbox}{milcVino}{milcGris}{Sentido de la evaluación}
Evaluar no es solo revisar una nota. Es reconocer qué aprendí, qué cambió en mí, qué emociones manejé, cómo me relaciono con la ciudad y la comunidad, y qué le dejaré a quien viene detrás.
\end{softbox}

\textbf{Mi reflexión liberadora en cinco dimensiones.} Completa cada frase iniciada:

\small
\begin{tabularx}{\textwidth}{>{\bfseries\RaggedRight\arraybackslash}p{3.4cm}Y>{\RaggedRight\arraybackslash}p{4.6cm}}
\rowcolor{milcVino}\color{white}Dimensión & \color{white}\bfseries Mi respuesta (frame starter) & \color{white}\bfseries Algo que descubrí\\
1. Personal & Lo que más cambió en mí fue \linead{6.5cm} & \linead{4.2cm}\\
2. Emocional & La emoción más fuerte fue \linead{2.5cm} y la regulé \linead{3cm} & \linead{4.2cm}\\
3. Ciudadana & En democracia esto importa porque \linead{6cm} & \linead{4.2cm}\\
4. Local (Cartago) & En mi barrio sería útil para \linead{6.5cm} & \linead{4.2cm}\\
5. Intergeneracional & A mi abuela/abuelo le diría \linead{6.5cm} & \linead{4.2cm}\\
\end{tabularx}
\normalsize

\begin{infoband}{milcVino}
\textbf{Para la próxima sustentación.} Vuelve a esta tabla en seis meses. Verás que las respuestas cambian — no porque lo aprendido cambie, sino porque tú cambias. La evaluación liberadora es una conversación contigo a través del tiempo.
\end{infoband}


\puente{Y para terminar, tres voces sobre la palabra que decide. Dussel sobre los informes que dan voz a propuestas necesarias, Marco Aurelio sobre la claridad como virtud comercial, Floridi sobre la propuesta documentada como ética profesional.}

\milcestacion{milcGradeDark}{Triángulo de pensamiento · cierre filosófico}

\begin{softbox}{milcVino}{milcGris}{Dussel · lente del nosotros}
\textit{Una propuesta documentada que sostiene voces necesarias es ya acto político; sin documento, la voz queda en queja.}

Cuando documentas una propuesta y la entregas al tomador de decisiones, estás transformando inconformidad en acción posible. La filosofía de la liberación pide ese paso: \textbf{convertir la queja en propuesta argumentada}. Sin documento, la voz se pierde en pasillos; con documento, queda registro que la institución debe responder.

\textbf{¿Mi informe transforma una necesidad real en propuesta argumentada, o solo es ejercicio escolar?}
\end{softbox}
\linead{\linewidth}\\[1mm]
\linead{\linewidth}

\begin{softbox}{milcMostaza}{milcGris}{Estoicismo · Marco Aurelio · cuidado interior}
\textit{La claridad del pedido es virtud comercial; esconder lo que se quiere es debilidad disfrazada de cortesía.}

Es tentador dar muchas vueltas antes del pedido (\emph{``esperamos que esta humilde propuesta pueda ser considerada en su valiosa agenda...''}). La sabiduría estoica recomienda lo opuesto: \emph{dí lo que pides con claridad en el primer párrafo}. La cortesía está en el tono general, no en esconder el pedido. La phronesis comercial consiste en \textbf{respetar el tiempo del lector siendo claro}.

\textbf{¿Mi resumen ejecutivo dice claramente lo que pido, o lo escondo por timidez?}
\end{softbox}
\linead{\linewidth}\\[1mm]
\linead{\linewidth}

\begin{softbox}{milcTurquesa}{milcGris}{Floridi · lente de la infoesfera}
\textit{La propuesta bien documentada con datos verificables es ética profesional en la era de la sobreabundancia de opinión.}

En la era de las redes sociales, hay muchas opiniones pero pocas propuestas documentadas. La ética profesional pide lo contrario: cuando propones algo, súbelo con datos, responsables, plazos. Tu informe comercial de hoy te entrena en una práctica que rige toda profesión que produzca propuestas reales: \textbf{argumentar con rigor, no solo con sentimiento}.

\textbf{¿Mi propuesta tiene rigor profesional, o es opinión disfrazada de informe?}
\end{softbox}
\linead{\linewidth}\\[1mm]
\linead{\linewidth}

\begin{infoband}{milcNegro}
\textbf{Nota para el docente.} Las tres formulaciones del triángulo son la síntesis del autor de la guía, no citas textuales de los pensadores. Si vas a citarlos en clase, remite a la obra original.
\end{infoband}

\milcestacion{milcVino}{Mi autoevaluación · marca una casilla por fila}

{\small
\setlength{\extrarowheight}{2.2pt}
\begin{tabularx}{\textwidth}{>{\RaggedRight\arraybackslash\bfseries}p{3.0cm}YYY}
\rowcolor{milcVino}\color{white}\bfseries Criterio & \color{white}\bfseries Lo logré & \color{white}\bfseries En proceso & \color{white}\bfseries Necesito apoyo\\[.6mm]
Comprensión del tema & \milcopcion{Explico las ideas con mis palabras.} & \milcopcion{Reconozco las ideas, falta precisión.} & \milcopcion{Necesito repasar lo básico.}\\[.8mm]
\rowcolor{milcGris}Pensamiento computacional & \milcopcion{Usé los pilares al armar mi producto.} & \milcopcion{Usé algunos pilares.} & \milcopcion{Necesito releer la estación 2.}\\[.8mm]
Aplicación y producto & \milcopcion{Mi producto cumple los criterios.} & \milcopcion{Está armado, con ajustes pendientes.} & \milcopcion{Aún está en construcción.}\\[.8mm]
\rowcolor{milcGris}Reflexión 5 dimensiones & \milcopcion{Respondí cada una con honestidad.} & \milcopcion{Respondí algunas con detalle.} & \milcopcion{Mis respuestas son muy generales.}\\[.8mm]
Triángulo de pensamiento & \milcopcion{Conecté las 3 voces con mi trabajo.} & \milcopcion{Conecté al menos una.} & \milcopcion{No las relaciono con mi proceso.}\\[.8mm]
\end{tabularx}}

\vspace{3mm}
\textbf{Hoy entendí que:}\\[1mm]
\linead{\linewidth}\\[1mm]
\linead{\linewidth}

\textbf{Mi compromiso para la próxima sesión es:}\\[1mm]
\linead{\linewidth}\\[1mm]
\linead{\linewidth}

\begin{softbox}{milcMostaza}{milcGris}{Cita de despedida · Marco Aurelio}
\textit{``El obstáculo en el camino se vuelve el camino. Lo que estorba al camino se vuelve camino.''}\\[1mm]
Cada error de hoy, bien observado, se convierte en pieza del camino que pisarás mañana.
\end{softbox}

\small
\begin{tabularx}{\textwidth}{>{\bfseries}p{3.6cm}Y}
\rowcolor{milcGradeDark!12}Nombre completo: & \linead{10cm}\\
Fecha: & \linead{4cm} \quad Curso/grupo: \linead{4cm}\\
Mi firma: & \linead{9cm}\\
\end{tabularx}
\normalsize

\begin{infoband}{milcGradeDark}
\textbf{Para la próxima sesión.} Llega con tu informe comercial completo y plan de entrega. En la sesión 8 vas a aprender a escribir \textbf{correos profesionales con IA como asistente}: el complemento natural del informe comercial cuando se envía por correo electrónico. El informe de hoy es el documento; el correo de mañana es la primera impresión.
\end{infoband}

\end{document}
